\section{Modus operandi}
\label{sec:modus-operandi}

Un primo aspetto del metodo di lavoro ha riguardato il confronto tra la conoscenza di dominio condivisa dai referenti aziendali e le evidenze emerse dai dati. Le indicazioni ricevute hanno guidato la comprensione del processo produttivo e la formulazione di ipotesi, poi verificate sul dataset disponibile. Quando sono emerse discrepanze, abbiamo approfondito l'analisi e chiesto chiarimenti, per comprenderne le cause e precisare le condizioni in cui le ipotesi risultavano valide.

Su questa base, nella preparazione dei dati abbiamo dato precedenza alla coerenza logica dei collegamenti tra le operazioni e dei passaggi di materiale, comprese diramazioni e ricongiungimenti. Abbiamo prima verificato il rispetto dei bilanci di massa e la compatibilità temporale delle lavorazioni. Quando i dati e la conoscenza di dominio lasciavano ambiguità, abbiamo introdotto ipotesi di lavoro per ottenere ricostruzioni plausibili, distinguendole dalle indicazioni aziendali e dai fatti osservati. Solo in seguito abbiamo svolto le analisi statistiche, evitando di considerare la sola rarità di un valore come motivo sufficiente per correggerlo.
